\subsection{Proof of Theorem~\ref{thm-simplest}}

Lemma~\ref{lem-invariance} implies that $\partial_I X_J$ is a polynomial in $X_0, X_1, \cdots, X_n$ with coefficients in $\FF_2$.  Consider the grading by $\ZZ^{n+1}$ on the  ring $\FF_2[\a{i, j}]$ such that $\a{i, j}$ has degree $e_j$. Here $e_0,\ldots, e_{n}$ is the standard basis for $\ZZ^{n+1}$. Let $\one = (1,\ldots, 1)$. Then $X_i$, $i=0,\ldots, n$, is homogeneous with
\[ \deg X_i = \one - e_i.\]
Since the vectors $\one - e_i$ are linearly independent, there can be at most one monomial $X_L$ in each degree. 
The partial derivative $\partial_{\a{i,j}}$ applied to a homogeneous polynomial reduces its degree by $e_j$ (or is zero). Since $X_J$ is homogeneous, so is $\partial_I X_J$, hence $\partial_I X_J$ is equal to a constant $c$ times a monomial $X_M$. Here $c\in \FF_2$, hence $\partial_I X_J$ is either $X_M$ or $0$.
Computing the degrees, the monomial $X_M$ must satisfy $X_I X_J = X_M X_{(0,1,\dots,n)}$.
Theorem~\ref{thm-simplest} has two cases depending on whether $\sqrt{X_M}$ exists or not.

\begin{lemma} 
% If $\sqrt{X_M}$ does not exist then $\partial_I X_J = 0$.
If $\partial_I X_J \neq 0$, then $\sqrt{X_M}$ exists.
\end{lemma}

\begin{proof}
Suppose that $\partial_I X_J = X_M$. By Lemma~\ref{lem-2deriv}, all partial derivatives  
 \[ \partial_{\a{i, j}} X_M = \partial_{\a{i, j}} \partial_I X_J\]
 vanish. This implies that $X_M$ is the square of a polynomial in $\FF_2[\a{i, j}]$. Since $X_M$ is the product of irreducible polynomials $X_i$, it follows that $X_M$ must be the square of a monomial $X_L$.
\end{proof}

\begin{lemma}
% If $\sqrt{X_M} = X_L$ then $\partial_I X_J = X_M$.
If $\sqrt{X_M}$ exists, then $\partial_I X_J \neq 0$.
\end{lemma}

\begin{proof}
We will prove that $\partial_I X_J \neq 0$ by induction on $n$. 
The base of the induction is $n=0$. In this case $\partial_I X_J = X_J =1$.

Consider now $n>0$.  Let $I=(i_1,\ldots, i_n)$ and $J=(j_1,\ldots,j_s)$, where $s$ is odd. By assumption, there exists a monomial $X_L$ such that $X_I X_J = X_L^2 X_{(0,1,\ldots,n)}$.  To prove that $\partial_I X_J \neq 0$, we may factor out squares in $X_J$ and assume that $J$ has no repeated entries.

There exists an entry in $J$, say $j_1$, such that $j_1 \neq i_r$ for $r=1,\ldots, n$. This follows from the fact that $X_{(0,1,\ldots,n)}$ of degree $n+1$ divides $X_I X_J$, but $X_I$ has degree $n$. Define the monomial  
  \[ \mu = \a{1,i_1} \a{1,j_1}^{s-1}.\]
Then the coefficient of $\mu$ in $X_J=X_{j_1} \cdots X_{j_s}$ is 
\[  Y_{j_1, i_1} Y_{j_2, j_1} \cdots Y_{j_s,j_1}.\]
Indeed, $X_{j_1}$ does not contain $\a{1,j_1}$. Hence $\a{1,j_1}^{s-1}$ must come from the factors $X_{j_2}, \ldots, X_{j_s}$  and $\a{1,i_1}$ from the factor $X_{j_1}$. As before, we have denoted by $Y_{i,j}$ the coefficient of $\a{1,j}$ in $X_i$.

Notice that the coefficient of $\a{1,j_1}^{s-1}$ in $\partial_I X_J$ is 
\[ \partial_{\a{2,i_2}} \partial_{\a{3,i_3}} \cdots \partial_{\a{n,i_n}} Y_{j_1, i_1} Y_{j_2, j_1} \cdots Y_{j_s,j_1}.\]
If this coefficient is nonzero then also $\partial_I X_J$ is nonzero. We claim that this coefficient being nonzero follows by induction from the case of dimension $n-1$.  From the matrix $(\a{i,j})$ we have removed row $1$ and column $j_1$, the derivative $\partial_I$ is replaced with $\partial_{I'}$, where $I' = (i_2,\ldots, i_n)$, and $X_J$ is replaced (using a similar notation in dimension $n-1$) with $X_{J'}$, where $J'=(i_1, j_2, \ldots, j_s)$. 

Let us check that we can apply the induction assumption to prove that $\partial_{I'} X_{J'} \neq 0$. The vectors $I'$ and $J'$ satisfy
\[ X_{I'} X_{J'} = \frac{X_I}{X_{i_1}} \cdot \frac{ X_{i_1} X_J }{ X_{j_1} } = \frac{X_L^2 X_{(0,1,\dots,n-1,n)}}{X_{j_1}}  \]
Since we assumed that $J$ does not contain repeated entries, $X_{j_1}$ appears in the numerator of the fraction to the first power. If we suppose that $j_1 = n$, then
\[ X_{I'} X_{J'} = X_L^2 X_{(0,1,\dots, n-1)}, \]
where $X_{j_1} = X_n$ does not appear in any monomial. By induction, $\partial_{I'} X_{J'} \neq 0$.
\end{proof}
